\section{Validação narrativa e comparação setorial}
\label{sec:eventos}

A avaliação sintética (Seção~\ref{sec:avaliacao}) é controlada porém artificial.
Falta o teste de realidade: \textbf{as anomalias detectadas coincidem com eventos
que de fato abalaram o mercado?}

\subsection{Cruzar com a história}
Montamos uma linha do tempo curada de eventos econômicos e políticos brasileiros
de 2010 a 2024 (crash da COVID, fraude da Americanas, Brumadinho, intervenção na
Petrobras, \emph{impeachment}, etc.). Cada evento é classificado como:
\begin{itemize}[noitemsep]
  \item \textbf{sistêmico}: afeta todo o mercado (ex.: COVID);
  \item \textbf{idiossincrático}: específico de uma empresa (ex.: fraude da
        Americanas $\to$ só AMER3).
\end{itemize}

\begin{definicao}[Tolerância de \emph{matching}]
Dizemos que um evento foi ``coberto'' se houve ao menos uma anomalia detectada
dentro de $\pm 30$ dias dele. A tolerância de 30 dias não é arbitrária: é o
próprio tamanho da janela (\texttt{window\_size}), a granularidade natural da
detecção. Datas de eventos também são âncoras aproximadas (o anúncio e a reação
do mercado podem diferir em dias).
\end{definicao}

\subsection{Os dois testes que isso permite}
\begin{itemize}[noitemsep]
  \item \textbf{Contágio (sistêmico):} um evento que abala tudo deve acender
        anomalia em \emph{todos} os ativos. A Figura~\ref{fig:teoria-m6} mostra o
        crash da COVID disparando nos quatro, de setores distintos.
  \item \textbf{Especificidade (idiossincrático):} um evento de uma empresa deve
        acender \emph{só} nela. A fraude da Americanas isola-se em AMER3;
        Brumadinho, em VALE3.
\end{itemize}

A combinação dos dois mostra que o modelo, treinado por ativo, captura tanto
choques amplos de mercado quanto problemas específicos --- o objetivo de
``generalização cruzada entre setores'' do projeto.

\begin{figure}[h]
\centering
\includegraphics[width=\textwidth]{m6_covid_contagio}
\caption{Crash da COVID (fev--mai/2020): os quatro ativos, de setores diferentes,
sinalizam anomalias concentradas no mesmo período --- evidência de contágio
sistêmico.}
\label{fig:teoria-m6}
\end{figure}

\begin{intuicao}[Por que cobertura ``só'' de 50\% não é ruim]
A cobertura agregada de eventos fica em torno de 50\%, mas o número cru engana:
o detector é calibrado para a \emph{cauda} (5\% maiores erros), então
deliberadamente ignora eventos de impacto difuso e dispara nos mais fortes. O que
importa é o \emph{casamento certo}: AMER3 cobre 9 de 10 dos seus eventos. Um
detector que acendesse em tudo teria 100\% de cobertura e nenhuma utilidade ---
de novo, o trade-off precision/recall da Seção~\ref{sec:avaliacao}.
\end{intuicao}

\begin{naprat}
Como há \emph{uma} ação por setor ($n=1$), a comparação setorial é um
\textbf{estudo de caso qualitativo}, não inferência estatística: diferenças
podem refletir a empresa, não o setor. Reportamos isso explicitamente para não
exagerar nas conclusões.
\end{naprat}
